\section{The smoothing criterion}\label{sec04:smoothing}

Throughout this section, $X$ is a connected normal projective complex
threefold with $\omega_X\cong\cO_X$ and $H^1(X,\cO_X)=0$, whose
singularities are isolated ordinary double points and the specified analytic
anticanonical cones over $\dP_5$, $\dP_6$, $\dP_7$, and
$\PP^1\times\PP^1$.  Fix a local component $S_p$
at every cone point, and write $S=(S_p)$.  We retain the entire nodal base
and the entire degree-five and $\PP^1\times\PP^1$ cone bases, and use the reduced
component at a degree-seven point.  The smooth crepant
resolution $Y_0\to X$, the link subspaces $R_{p,S_p}$, and their relation
space
\begin{equation}\label{eq04:relations}
 K_S=\ker\left(\bigoplus_p R_{p,S_p}
                   \longrightarrow H_2(Y_0;\QQ)\right)
\end{equation}
are as in Section~\ref{sec:deformation}.  All tangent identifications below
use contraction with one fixed generator of $H^0(X,\omega_X)$.  In
particular, coefficients in~\eqref{eq04:relations} are the resulting
homological coordinates; they are not independently normalized local
equation parameters.

Choose an analytic semiuniversal deformation of $X$, analytic local
semiuniversal deformations of its singularities, and analytic classifying
maps from the first base to the latter bases.  We denote by $\Def^S(X)$
the scheme-theoretic inverse image of the product of the selected local
components.  These choices specify the analytic subgerm and its completed
local ring.  We use $\mathbf{Def}_X$ for the functor of marked deformations
over local Artin $\CC$-algebras, and $\Def(X)$ for an analytic
semiuniversal base.

\subsection{Global partial resolutions}

\begin{proposition}\label{prop04:globalmodel}
The local morphisms of Proposition~\ref{prop03:localmodels}, together
with the identity near degree-five and original
$\PP^1\times\PP^1$ cones, determine a
proper crepant analytic modification $\pi:Z\to X$.  Its singularities are
the original nodes, two or three nodes over each degree-six point according
to its selected component, and one anticanonical
$\PP^1\times\PP^1$ cone over each degree-seven point, together with
the original degree-five and $\PP^1\times\PP^1$ cones.  Moreover,
\begin{equation}\label{eq04:cy}
 R\pi_*\cO_Z\simeq\cO_X,\qquad
 \omega_Z\cong\cO_Z,\qquad H^1(Z,\cO_Z)=0.
\end{equation}
Every smooth resolution of $Z$ is Moishezon and satisfies the
$\partial\bar\partial$-lemma.
\end{proposition}

\begin{proof}
Choose mutually disjoint sufficiently small neighborhoods $U_p$ of the
degree-six and degree-seven cone points, and restrict the local morphisms to proper morphisms
$W_p\to U_p$.  Glue the $W_p$ to the unchanged complement of these
degree-six and degree-seven points using their specified isomorphisms
over $U_p\setminus\{p\}$.  This complement includes every original
degree-five and $\PP^1\times\PP^1$ cone.
The images of these charts are open.  Two points with different images in
$X$ are separated by inverse images of disjoint neighborhoods in $X$;
two distinct points over the same cone point are separated in the
Hausdorff space $W_p$.  Thus the glued space is Hausdorff.  The map is
proper over every $U_p$ and over the complement of these points,
and properness is local on the target.  This proves properness globally,
and hence compactness of $Z$.

The local models are Gorenstein and canonical.  In particular, $X$ and
$Z$ have rational singularities.  If $g:\widetilde Z\to Z$ is a
resolution, then $\pi g$ is a resolution of $X$ and
$Rg_*\cO_{\widetilde Z}\simeq\cO_Z$,
$R(\pi g)_*\cO_{\widetilde Z}\simeq\cO_X$.  Composition gives the
first assertion of~\eqref{eq04:cy}, and Leray gives its last assertion.
Local crepancy means that the pullback of the chosen volume form extends
to a generator of the dualizing sheaf on every local model.  These
generators agree on the overlaps, giving $\omega_Z\cong\cO_Z$.

A smooth resolution is compact and bimeromorphic to projective $X$, hence
Moishezon.  A smooth projective modification of it exists, and the
$\partial\bar\partial$-lemma descends under a proper bimeromorphic map
between compact complex manifolds
\cite[Theorem~5.22 and Corollary~5.23]{dgms1975}.
\end{proof}

Let $Y\to Z$ be the smooth crepant resolution obtained by resolving the
new nodes on the sides specified in Proposition~\ref{prop03:localmodels}
and resolving the degree-five and quadric cones by their vertex blowups.
At a degree-six
point this recovers the original neighborhood in $\operatorname{Tot}K_E$.
At a degree-seven point it differs from $Y_0$ by the described flop.
Near a degree-five or original $\PP^1\times\PP^1$ cone the two
resolutions agree.
Removing a smooth curve from a smooth threefold does not change $H_2$:
the relative groups in degrees two and three vanish by the real
codimension-four Thom isomorphism.  The common complements therefore
identify $H_2(Y_0;\QQ)$ and $H_2(Y;\QQ)$.  The degree-seven root is
represented on that complement by the difference of the two ruling
classes, so the identifications preserve the local substitutions used
in~\eqref{eq04:relations}.

\subsection{Blowing down Artin deformations analytically}

The following lemma supplies a morphism of deformation functors even
when $Z$ is not algebraic.  Its sheaf-theoretic part is the construction
in \cite[Proposition~2.12 and Corollary~2.13]{sano2017}; we include the
analytic realization.

\begin{lemma}[analytic blow-down]\label{lem04:blowdown}
Let $f:Z\to X$ be a proper holomorphic morphism of complex spaces with
$f_*\cO_Z=\cO_X$ and $R^1f_*\cO_Z=0$.  A flat deformation $Z_A$ over
a local Artin $\CC$-algebra $A$ determines a flat analytic deformation
\[
 X_A=(|X|,f_*\cO_{Z_A})
\]
and a proper morphism $f_A:Z_A\to X_A$ deforming $f$.  This construction
is natural in marked deformations and commutes with every Artin base
change.
\end{lemma}

\begin{proof}
The marking identifies the underlying topological space of $Z_A$ with
$|Z|$, so the displayed pushforward is defined as a sheaf of $A$-algebras.
Let $U\subset X$ be a sufficiently small Stein open subset and put
$V=f^{-1}(U)$, $\mathcal F=\cO_{Z_A}|_V$.  Leray and Cartan's theorem~B
give $H^1(V,\cO_Z)=0$.  For every finite $A$-module $M$ we claim
\begin{equation}\label{eq04:modulevanish}
 H^1(V,\mathcal F\otimes_A M)=0.
\end{equation}
Indeed, a composition series of $M$ has successive quotients $\CC$.
Flatness of $\mathcal F$ preserves its short exact sequences under
tensor product.  The long cohomology sequences, starting from
$\mathcal F\otimes_A\CC=\cO_Z|_V$, prove
\eqref{eq04:modulevanish} by induction.

It follows that $M\mapsto H^0(V,\mathcal F\otimes_A M)$ is exact on
finite $A$-modules.  Applying it to a finite presentation of $M$ by free
modules gives the natural isomorphism
\begin{equation}\label{eq04:modulebc}
 H^0(V,\mathcal F)\otimes_A M
       \xrightarrow{\ \sim\ }H^0(V,\mathcal F\otimes_A M).
\end{equation}
Consequently $H^0(V,\mathcal F)$ is $A$-flat: exactness on finite
modules suffices over the Noetherian ring $A$.  Passing to stalks on
$X$ proves flatness of $f_*\cO_{Z_A}$ and the corresponding sheaf
version of~\eqref{eq04:modulebc}.  In particular its reduction modulo
the maximal ideal $\mathfrak m_A$ is $\cO_X$.

To prove analyticity, take $U$ small enough to be a closed analytic
subspace of a polydisc $D$.  Equation~\eqref{eq04:modulebc}, with
$M=\CC$, lifts its embedding coordinates to global sections of
$\mathcal F$.  Together with the structure morphism to
$\operatorname{Spec}_{\mathrm{an}}A$ these sections define a holomorphic
map
\[
 g_A:V_A\longrightarrow D\times\operatorname{Spec}_{\mathrm{an}}A.
\]
Its underlying continuous map is the proper composite
$V\to U\hookrightarrow D$, hence $g_A$ is proper.  Grauert's direct
image theorem for complex spaces, including nonreduced spaces,
implies that $g_{A*}\cO_{V_A}$ is coherent on the polydisc thickening
\cite[\S6, Hauptsatz~I]{grauert1960}.  The natural algebra map
\[
 \cO_{D\times\operatorname{Spec}_{\mathrm{an}}A}
                   \longrightarrow g_{A*}\cO_{V_A}
\]
is surjective modulo $\mathfrak m_A$, since its reduction is
$\cO_D\to\cO_U$.  Its coherent cokernel $\mathcal C$ therefore
satisfies $\mathcal C=\mathfrak m_A\mathcal C$, and nilpotence gives
$\mathcal C=0$.  Its kernel is a coherent ideal.  The resulting closed
analytic subspace has underlying ringed space
$(U,f_*\cO_{Z_A}|_U)$.

This ringed space is independent of the lifted coordinates.  Thus the
local analytic realizations glue, and the adjunction counit induces
$f_A$.  Locally this is the holomorphic map to the closed subspace just
constructed; its central restriction is $f$ and its underlying map is
proper.  Finally, if $A\to A'$ is any homomorphism of local Artin
$\CC$-algebras, then $A'$ is a finite $A$-module.  The sheaf version
of~\eqref{eq04:modulebc} identifies the blow-down of
$Z_A\times_A A'$ with $X_A\times_A A'$.  These natural tensor
identifications respect composition.  An isomorphism of marked source
deformations pushes forward to an isomorphism of the constructed target
deformations, proving the functorial assertion.
\end{proof}

\subsection{Unobstructedness of the partial resolution}

\begin{lemma}\label{lem04:unobstructed}
The functor $\mathbf{Def}_Z$ is unobstructed.  Both
$\mathbf{Def}_Z$ and $\mathbf{Def}_X$ are prorepresentable, and
Lemma~\ref{lem04:blowdown} induces a morphism of their formal bases
\begin{equation}\label{eq04:beta}
 \beta:\widehat{\Def}(Z)\longrightarrow\widehat{\Def}(X).
\end{equation}
\end{lemma}

\begin{proof}
By Proposition~\ref{prop04:globalmodel}, $Z$ is compact, normal and
Gorenstein with trivial dualizing sheaf, $H^1(Z,\cO_Z)=0$, and isolated
rational, hence Du Bois, singularities.  A resolution satisfies the
$\partial\bar\partial$-lemma.  Friedman's theorem
\cite[Theorem~3.10]{friedman2026} gives the generalized $T^1$ lifting
property.  We explain why here it gives ordinary $T^1$ lifting.

Set $A_k=\CC[t]/(t^{k+1})$ and
$B_k=A_k[\epsilon]/(\epsilon^2)$.  For a deformation $Z_k/A_k$ let
$Z_{k-1}$ denote its restriction.  Relative first-order deformations,
with an identification over $\epsilon=0$, form $T^1(Z_k/A_k)$.
There is an exact sequence and its local counterpart
\[
 \begin{aligned}
 T^1(Z_k/A_k)&\longrightarrow T^1(Z_{k-1}/A_{k-1})
       \xrightarrow{\delta}\mathbb T^2_Z,\\
 T^1_{\mathrm{loc}}(Z_k/A_k)&\longrightarrow
 T^1_{\mathrm{loc}}(Z_{k-1}/A_{k-1})
       \xrightarrow{\delta_{\mathrm{loc}}}
       H^0(Z,\mathcal T^2_Z).
 \end{aligned}
\]
Here $\mathbb T^2_Z=\operatorname{Ext}^2(\Omega^1_Z,\cO_Z)$ and
$\mathcal T^2_Z=\mathcal{E}xt^2(\Omega^1_Z,\cO_Z)$ are the global
and local obstruction spaces used in the lifting theorem.
Localization $\ell:\mathbb T^2_Z\to H^0(Z,\mathcal T^2_Z)$
commutes with the connecting maps.  Generalized lifting means that
$\ell$ is injective on $\operatorname{im}\delta$ for every $k$ and
every $Z_k$.

The full local deformation bases of $Z$ are smooth: this is the usual
nodal calculation, Proposition~\ref{prop03:degree5family}, and
Lemma~\ref{lem03:quadric}.  This implies
surjectivity of the displayed local $T^1$ map with its markings, as
follows.  For a fixed local deformation over $A_k$ and a marked relative
deformation over $B_{k-1}$, use the specified identification over
$A_{k-1}$ to glue them over
\[
 D_k=A_k\times_{A_{k-1}}B_{k-1}.
\]
The deformation category of analytic spaces has the
Rim--Schlessinger gluing property: the local structure sheaves are
formed by the fibre product over the common restriction, and the
result comes with compatible identifications on both factors.  This
is gluing with a specified identification, not merely a choice of
isomorphism classes; see also the categorical formulation in
\cite[\S1]{Gross97rev}.  The natural map
$B_k\to D_k$ is surjective with kernel $(\epsilon t^k)$.
Smoothness of the local deformation functor lifts the glued object to
$B_k$.  Choose an isomorphism from its restriction to $D_k$ to that
object, and restrict this same isomorphism to the two factors.
The resulting identifications agree over $A_{k-1}$ and give exactly
the required relative first-order lift.  Infinitesimal automorphisms
can change these choices but do not prevent their existence; local
prorepresentability is not used.

Thus $\delta_{\mathrm{loc}}=0$, so $\ell\delta=0$ and injectivity
on $\operatorname{im}\delta$ gives $\delta=0$.  Ordinary $T^1$
lifting proves unobstructedness
\cite[Theorem~3.3]{friedman2026}.  In particular this argument concerns
the full base, including its nilpotent structure; it does not infer
smoothness from a statement about reduced supports.

For completeness, the vanishing of infinitesimal automorphisms follows
for both $T=X$ and $T=Z$ from the same argument.  Choose an equivariant
resolution $h:\widetilde T\to T$.  Canonical singularities allow the
volume form to pull back holomorphically, and contraction gives an
injection $\Theta_{\widetilde T}\to\Omega^2_{\widetilde T}$.
Rationality, Hodge symmetry on the Moishezon resolution, and Serre
duality give
\[
 h^0(\widetilde T,\Omega^2_{\widetilde T})
 =h^2(\widetilde T,\cO_{\widetilde T})
 =h^2(T,\cO_T)=h^1(T,\cO_T)=0.
\]
Since $h_*\Theta_{\widetilde T}=\Theta_T$, this proves
$H^0(T,\Theta_T)=0$, and the deformation functors are
prorepresentable; compare \cite[Lemma~3.2]{friedman2026}.
The natural Artin transformation supplied by
Lemma~\ref{lem04:blowdown} therefore gives~\eqref{eq04:beta}.
\end{proof}

\subsection{Lifting the entire selected tangent space}

\begin{lemma}\label{lem04:globaltangent}
Put $V_S=T_0\Def^S(X)$.  There is an exact sequence
\begin{equation}\label{eq04:tangent}
 0\longrightarrow H^1(X,\Theta_X)\longrightarrow V_S
       \longrightarrow K_S\otimes_{\QQ}\CC\longrightarrow0,
\end{equation}
and every vector in $V_S$ belongs to the image of $d\beta$.
\end{lemma}

\begin{proof}
The exact sequence is the restriction of Theorem~\ref{thm:D} to the
selected local tangent subspaces.  We construct lifts through $d\beta$,
including lifts of its locally trivial kernel.

Let $\kappa\in K_S\otimes\CC$.  At a degree-six $A_1$ point set
$C_0=e_1$, $C_3=h-e_2-e_3$.  Since the convention of
Section~\ref{sec:deformation} is $\ell=C_0-C_3$, replace the
coefficient $\lambda\ell$ by nodal coefficients
$(\lambda,-\lambda)$ at $(C_0,C_3)$.  At an $A_2$ point replace
$\sum_i\mu_i e_i$, $\sum_i\mu_i=0$, by the three nodal coefficients
$\mu_i$.  At a degree-seven point use the same multiple of the
quadric ruling difference.  Retain the old nodal coefficients and
the ruling-difference coefficients at original $\PP^1\times\PP^1$
cones, and retain the whole local vector at each degree-five point.  The
curve identifications of Proposition~\ref{prop03:localmodels} and the
common complements identify the sum of this tuple $\nu$ with the
sum of $\kappa$ in $H_2(Y;\CC)$; it is therefore zero.

Write $F$ for the exceptional locus of $Y\to Z$, with local pieces
$F_q\subset Y_q$.  The local Hodge comparisons identify the components
of $\nu$ with local tangent classes and with their images in
$H^2_{F_q}(Y_q,\Omega^2_{Y_q})$.  In the support sequence
\begin{equation}\label{eq04:support}
 H^1(Y\setminus F,\Omega^2)
 \longrightarrow\bigoplus_q H^2_{F_q}(Y_q,\Omega^2)
 \longrightarrow H^2(Y,\Omega^2_Y),
\end{equation}
their image in the last group is the class of the corresponding
curves.  Under the injection
$H^2(Y,\Omega^2_Y)\hookrightarrow H^4(Y;\CC)$ supplied by the
$\partial\bar\partial$-lemma it is the Poincar\'e dual of their
homological sum, hence zero.  Exactness gives a class on
$Y\setminus F$.  Contraction with the volume form and the
depth-three comparison of deformations with the regular locus give
a tangent $\eta\in\mathbb T^1_Z$ with local tuple $\nu$, as in the
proof of Theorem~\ref{thm:D}.

Let $E'$ be the full exceptional set of $Y\to X$.  Restriction from
$Y\setminus F$ to $Y\setminus E'$ commutes with the boundary maps
in~\eqref{eq04:support} and the inclusion of supports
$H^2_F(Y,\Omega^2_Y)\to H^2_{E'}(Y,\Omega^2_Y)$.  At a degree-six
point this sends the nodal tuple to its sum of curve classes in $E$;
at a degree-seven point it sends the ruling difference to the original
root; at a degree-five or original $\PP^1\times\PP^1$ cone
it is the identity.
These are exactly the local components of $\kappa$.
The construction of blow-down agrees with restriction on the common
regular locus.  Thus $d\beta(\eta)$ has the prescribed local
tangent on $X$, including at points where its coefficient is zero.

Given any desired vector $v\in V_S$ over $\kappa$, its difference
from $d\beta(\eta)$ belongs to $H^1(X,\Theta_X)$.  Choose a Stein
cover $\{U_i\}$ with exactly one member containing each singular
point and all other members avoiding the singular set.  Every overlap
lies in the locus where $\pi$ is an isomorphism.  Represent the
difference by a \v{C}ech cocycle $\{\xi_{ij}\}$ of vector fields.
Glue the trivial deformations of $\pi^{-1}(U_i)$ by the
automorphisms $1+\epsilon\xi_{ij}$ on the overlaps.  The additive
cocycle condition gives the gluing identity modulo $\epsilon^2$.
Pushing forward these trivial local deformations and their transitions
recovers the original cocycle on $X$.  Hence this locally trivial
global tangent also lifts, and adding its lift to $\eta$ gives $v$.
\end{proof}

\Needspace{7\baselineskip}
\subsection{Smoothness of a selected deformation base}

\begin{theorem}\label{thm04:smoothbase}
Suppose the projection $K_S\to R_{p,S_p}$ is nonzero at every
degree-six point $p$.  Then $\Def^S(X)$ is smooth, of dimension
\[
 h^1(X,\Theta_X)+\dim_{\QQ}K_S.
\]
Every tangent in $T_0\Def^S(X)$ is realized by a convergent curve in
this germ.  No nonzero-projection assumption is required at nodes,
degree-five or degree-seven points, or $\PP^1\times\PP^1$ cones.
\end{theorem}

\begin{proof}
Let $B$ be the completed local ring of $\Def(Z)$.  By
Lemma~\ref{lem04:unobstructed} it is a formal power-series ring, in
particular a domain.  At a degree-seven point the local base is
$\CC[[s,v]]/(v^2,sv)$, so the pullback of $v$ to $B$ vanishes.
At a degree-six point the local base is
$\CC[[s,u,v]]/(su,sv)$, with $A_1$ component $u=v=0$ and
$A_2$ component $s=0$.  If $A_1$ was selected, the nonzero
projection hypothesis and Lemma~\ref{lem04:globaltangent} imply
that the pullback of $s$ has a nonzero linear term.  It is nonzero in
$B$, and the two equations force the pullbacks of $u,v$ to vanish.
For a selected $A_2$ component, one of $u,v$ has a nonzero linear
term, forcing the pullback of $s$ to vanish.  Thus $\beta$ factors
through the full selected formal subspace.  Nodes, degree-five
cones, and original $\PP^1\times\PP^1$ cones impose no branch equations, since their
entire local bases are retained.  Its differential onto
that subspace is surjective by Lemma~\ref{lem04:globaltangent}.

Write its completed target ring minimally as
$A=\CC[[x_1,\ldots,x_r]]/I$, with $I\subset(x_1,\ldots,x_r)^2$.
Surjectivity of the differential means that the images of the $x_i$
in $B$ have independent linear terms.  Extend them to a formal
coordinate system of $B$.  The composite
$\CC[[x_1,\ldots,x_r]]\to B$ is injective, whereas every element
of $I$ maps to zero.  Hence $I=0$.  Equation~\eqref{eq04:tangent}
gives the dimension.  An analytic local ring is regular if its
completion is regular, so the analytic subgerm $\Def^S(X)$ is
smooth.  Analytic coordinates on that germ realize every tangent by
a convergent curve.  This last step does not require convergence of
the formal map $\beta$.
\end{proof}

\subsection{Necessity from local periods}

We recall the form of the period argument that will be used after a
simultaneous partial resolution.  Its hypotheses require compactness;
the intermediate space need not be projective.

\begin{lemma}[period necessity]\label{lem04:periodnecessity}
Let $T$ be a connected compact normal Gorenstein threefold with
$\omega_T\cong\cO_T$, $H^1(T,\cO_T)=0$, and a smooth crepant
resolution $\widetilde T\to T$.  Suppose its singularities are nodes
and anticanonical cones over $\dP_5$, $\dP_6$, $\dP_7$, or
$\PP^1\times\PP^1$, and that it has a
convergent proper smoothing which uses the $A_1$ component at every
degree-six point.  Then the relation space of its selected local
subspaces in $H_2(\widetilde T;\QQ)$ contains a vector nonzero
on every rank-one summand and satisfying
$\prod_{i=1}^5\lambda_i\ne0$ at every degree-five point.
\end{lemma}

\begin{proof}
We give the argument in a form applicable to the possibly nonprojective
space $T$; the local period computations and link markings are those
of \cite[\S\S3--4]{PaperIV}, with the
$\PP^1\times\PP^1$ case supplied below and the degree-five
case supplied by Lemma~\ref{lem03:degree5periods}.
Write $T_t$ for a smooth nearby fibre.
Remove disjoint small singular neighborhoods and let $W$ be the
common exterior with boundary $\bigsqcup_pL_p$.  Smooth local
triviality along $W$ gives decompositions
\[
 T_t=W\cup\bigsqcup_p M_p,\qquad
 \widetilde T=W\cup\bigsqcup_p N_p,
\]
where $M_p$ is the local Milnor fibre and $N_p$ a resolution
neighborhood, with the canonical common link identifications.

The smoothing has a relative holomorphic volume form after shrinking
the disc.  Here are the details needed without a polarization.  The
relative dualizing sheaf is invertible because the family is flat
with Gorenstein central fibre, and is compatible with base change.
Its first Chern class restricts to zero on $W$, by smooth triviality
and $\omega_T\cong\cO_T$, and on each $M_p$, by a local
trivialization near the central singularity.  The links satisfy
$H^1(L_p;\ZZ)=0$: for nodes this follows from
$L_p\cong S^2\times S^3$, and for the cone links it follows from
the homotopy sequence of the circle bundle with nonzero Euler class
$K_E$ over the simply connected del Pezzo surface.  Mayer--Vietoris
therefore injects $H^2(T_t;\ZZ)$ into the sum of the corresponding
groups of the pieces, so $c_1(\omega_{T_t})=0$.  Semicontinuity gives
$H^1(T_t,\cO_{T_t})=0$; the exponential sequence implies
$\omega_{T_t}\cong\cO_{T_t}$.  Connectedness and Grauert base
change now give a relative section $\Omega$ generating the canonical
sheaf on every fibre.  Its restriction to a smooth fibre is closed.

At a node with local equation $xy-zw=b(t)$, the period of the
residue form on the vanishing sphere is $c\,b(t)$, where $c\ne0$.
Indeed, after writing the quadratic form as $\sum_i z_i^2$, the
sphere is $\sqrt b\,S^3$ and the residue form scales by $b$;
its period at $b=1$ is $\pi^2$ up to orientation and the linear
change of coordinates.  The global form is a holomorphic unit times
this residue form.  Uniformly on the shrinking sphere its period is
therefore
\[
 b(t)\bigl(c\,g(p,0)+o(1)\bigr),\qquad g(p,0)\ne0,
\]
which is nonzero for every sufficiently small $t\ne0$, regardless
of the vanishing order of $b(t)$.

At a $\PP^1\times\PP^1$ cone, Lemma~\ref{lem03:quadric}
identifies the local smoothing with a pullback of
\eqref{eq03:quadfamily}, say with parameter $b(t)$.
The nodal residue form on its double cover is invariant under the
involution.  The total quotient $\AAA^4/\boldsymbol\mu_2$ is
Gorenstein because $-I\in\operatorname{SL}_4$; adjunction for its
parameter function shows that the descended residue generates the
relative dualizing sheaf, including at the central vertex.
This remains true after pullback along $b(t)$.
On the noncentral fibre the covering
is free, and the antipodal involution preserves the orientation of
the vanishing $3$-sphere.  Thus the period on the quotient cycle
$\sqrt{b(t)}\,S^3/\boldsymbol\mu_2$ is half its period on the
sphere.  The global volume form differs from the descended residue
form by a holomorphic unit $g$, and the cycle contracts to the
central vertex.  Consequently its period is
\[
 b(t)\bigl(\tfrac12 c\,g(p,0)+o(1)\bigr),\qquad c\ne0.
\]
It is nonzero for all sufficiently small $t\ne0$, for any
vanishing order of $b(t)$.  Lemma~\ref{lem03:quadricmilnor}
identifies this cycle with a generator of $H_3(M_p;\QQ)$ and
identifies the local boundary map with the ruling-difference line.
The link has $H^1(L_p;\ZZ)=0$ even though its fundamental group
is $\ZZ/2$, as required in the preceding canonical-bundle argument.

For the degree-six and degree-seven branches in the lemma, the smooth affine fibre is
$M=V\setminus E$, with $V=(\PP^1)^3$ for degree six and
$V=\operatorname{Bl}_{\mathrm{pt}}\PP^3$ for degree seven, and
$E\in|-\tfrac12K_V|$.  The comparison with a compact local Milnor
fibre, including its boundary marking, is
\cite[Lem.~4.2]{PaperIV}.  Weak Lefschetz makes
$H^2(V;\QQ)\to H^2(E;\QQ)$ injective; the respective Picard
ranks are $(3,4)$ and $(2,3)$.  Since $H^3(V;\QQ)=0$, the pair
sequence gives $b_3(M)=1$.  More explicitly, take the character
$w=\chi^m$ with $m=(1,1,1)$ on $(\PP^1)^3$, and $m=(1,1,-1)$
on the blowup fan with rays $e_1,e_2,e_3,-e_1-e_2-e_3,e_1+e_2+e_3$.
The closure of $w=1$ is $E$, and
\[
 \Omega_M=\frac{w}{(w-1)^2}
      \frac{dx}{x}\wedge\frac{dy}{y}\wedge\frac{dz}{z}
\]
has order zero at each toric boundary divisor: the order is
$\langle m,v\rangle-2\min(\langle m,v\rangle,0)-1=0$, since
$\langle m,v\rangle=\pm1$.  Its only pole is the double pole at
$E$, so it generates $\omega_M$.  In adjacent toric charts it is
$\pm dp\wedge dq\wedge dz/(q-pz)^2$.  The transitions are
\[
 (p',q',z')=(pz^c,qz^c,z^{-1}),\qquad
 c=0\ \text{or}\ -1
\]
in the two respective cases.  Fix $\varepsilon>0$ and $R>1$.
A closed three-cycle consists of the two caps
\[
 \{p=0,\ q=\varepsilon R e^{i\psi},\ |z|\le R\},\qquad
 \{p'=\varepsilon R^c e^{i\psi},\ q'=0,\ |z'|\le R^{-1}\}
\]
and the intervening chain
\[
 p=r\varepsilon e^{i\theta},\quad
 q=-(1-r)\varepsilon R e^{i(\theta+\varphi)},\quad
 z=R e^{i\varphi},\qquad 0\le r\le1.
\]
The boundaries agree under the stated transition, with opposite
orientations.  On the chain $q-pz=-\varepsilon R
e^{i(\theta+\varphi)}$, so it misses $E$; the caps also miss $E$.
The form vanishes on the caps and pulls back on the chain to
$\pm dr\wedge d\theta\wedge d\varphi$.  Its period is consequently
$\pm4\pi^2$, in particular nonzero.

The one-dimensional local smoothing is equivariant with positive
weights on the cone coordinates and weight one on its parameter
$b$.  Graded Nakayama gives a homogeneous generator of its relative
canonical sheaf near the vertex.  Its nonvanishing locus is invariant
and contains the vertex, hence contains the entire fibre $b=1$:
every orbit in that fibre contracts to the vertex.  Its restriction
is therefore a constant multiple of the preceding form.
The multiplier is constant because an algebraic unit on $V\setminus E$
has divisor supported on $E$, and $\operatorname{Pic}(V)$ is
torsion-free with $[E]\ne0$.  Scaling the fixed compact cycle into
the fibre over $b(t)$ contracts it to the vertex.  If the homogeneous
form has weight $a$, its global period is
\[
 b(t)^a\bigl(c\,g(p,0)+o(1)\bigr),\qquad c\ne0.
\]
It is again nonzero for small $t\ne0$.
At a degree-five point, Lemma~\ref{lem03:degree5periods}
provides five classes $\nu_{p,i}$, the images of its boundary
classes $\beta_i$, with nonzero periods.  Their dual link
functionals are the five $\lambda_i$.
Every one of these specified local cycles survives in
$H_3(T_t;\QQ)$: otherwise it would bound a rational chain and
have zero period against the closed form $\Omega_t$.

By Poincar\'e duality choose, for each specified local cycle, a
rational class $\gamma\in H_3(T_t;\QQ)$ pairing nontrivially
with it.
Its Mayer--Vietoris boundary belongs to
\[
 \ker\left\{\bigoplus_q H_2(L_q;\QQ)
 \longrightarrow H_2(W;\QQ)\oplus\bigoplus_qH_2(M_q;\QQ)\right\}.
\]
It therefore gives a relation in $H_2(\widetilde T;\QQ)$ and
belongs at each point to the selected local subspace.
Lefschetz duality identifies the local restriction of $\gamma$
with its pairing against $H_3(M_p;\QQ)$.  At a rank-one germ,
the boundary map
$H_3(M_p,L_p;\QQ)\to H_2(L_p;\QQ)$ is an isomorphism onto the
root line: its image is that line by the local Milnor-fibre
calculation, and its domain has dimension one.  The corresponding
coefficient is therefore nonzero.  At a degree-five point,
\eqref{eq03:degree5boundary} is an isomorphism onto the full
link space.  Compatibility of the boundary with the Lefschetz
pairing identifies its pairing with $\beta_i$ with the preceding
nonzero pairing against $\nu_{p,i}$.  By
Lemma~\ref{lem03:degree5periods} this is a nonzero multiple of
$\lambda_i$ on the resulting local relation vector.
Thus each required scalar functional is nonzero on some rational
global relation.  Avoiding the finitely many proper kernels gives
one relation on which all of them are nonzero.
\end{proof}

\subsection{The exact criterion}

\begin{theorem}\label{thm04:criterion}
The threefold $X$ has a convergent one-parameter smoothing inducing
$S$ if and only if $K_S$ contains a vector nonzero on every
rank-one summand, satisfying $\prod_{i=1}^5\lambda_i(\kappa_p)\ne0$
at every degree-five point, and, at every $A_2$-assigned point, with
\[
 \kappa_p=\mu_1e_1+\mu_2e_2+\mu_3e_3,\qquad
 \mu_1+\mu_2+\mu_3=0,\qquad \mu_1\mu_2\mu_3\ne0.
\]
Equivalently, none of these individual scalar functionals vanishes
identically on $K_S$.
\end{theorem}

\begin{proof}
Suppose first that such a relation exists.  It gives a vector of
$V_S$ by~\eqref{eq04:tangent} and implies the nonzero-projection
hypothesis of Theorem~\ref{thm04:smoothbase}.  Realize that tangent
by a convergent curve in $\Def^S(X)$.  At each rank-one point the
local smoothing parameter has nonzero derivative.  At each $A_2$
point all three local discriminant factors have nonzero derivatives,
by the normalized identifications of
Proposition~\ref{prop03:localmodels}.  At each degree-five point
the five discriminant factors have nonzero derivatives by
Lemma~\ref{lem03:degree5marking}.  Thus the local fibres are
smooth for every sufficiently small nonzero parameter.  Properness
and openness of smoothness exclude additional singularities outside
neighborhoods of the finitely many original singularities.  The
curve therefore gives a smoothing of $X$.

Conversely, let $\mathcal X\to\Delta$ be any convergent proper
smoothing inducing $S$.  At each $A_2$ point pull back the
simultaneous incidence modification of
Proposition~\ref{prop03:a2family} along its actual analytic
classifying arc.  Since this arc is a smoothing, it meets the
central point of the local base only at $t=0$, after shrinking.
The pulled-back modification is proper and flat over $\Delta$,
is an isomorphism off the central cone point, and is an isomorphism
over every noncentral fibre.  Glue these morphisms to the identity
on the complement of those central points in $\mathcal X$.  The
Hausdorffness and properness argument of
Proposition~\ref{prop04:globalmodel} applies to this cover of the
total space; flatness is local on the source.  We obtain a proper
flat family $\mathcal T\to\Delta$ with $T_t\cong X_t$ for
$t\ne0$.

Its central fibre $T$ replaces each $A_2$ point by three nodes and
leaves all other germs unchanged.  Its smooth crepant resolution is
$Y_0$: the three new exceptional curves are $e_1,e_2,e_3$ in the
original del Pezzo surface.  Rationality and crepancy, as in
\eqref{eq04:cy}, give $\omega_T\cong\cO_T$ and
$H^1(T,\cO_T)=0$.  The remaining degree-six points use $A_1$.
Lemma~\ref{lem04:periodnecessity} therefore supplies a relation
nonzero at every node and every remaining rank-one cone subspace,
and with all five $\lambda_i$ nonzero at each degree-five point.  In
particular, the three new nodes over any one $A_2$ point have
nonzero coefficients $\mu_1,\mu_2,\mu_3$.

Pair this relation with the original exceptional divisor $E$ in
$Y_0$.  Every term belonging to another original singular point has
zero intersection with $E$, whereas
$E\cdot e_i=K_E\cdot e_i=-1$.  Hence
$\mu_1+\mu_2+\mu_3=0$.  Replacing the three nodal terms by
$\sum_i\mu_i e_i$ reconstructs a vector in $K_S$ with all the
required nonvanishing conditions.  This argument uses the entire
classifying arc; it imposes no assumption on its first nonzero order
and excludes an escape from a discriminant line at any higher order.
Finally, a finite-dimensional rational vector space is not a union
of finitely many proper linear subspaces, which proves the last
equivalence.
\end{proof}

In particular, when no degree-six points occur, every selected base
is smooth, even if a coordinate at a node, a degree-seven point,
or a $\PP^1\times\PP^1$ cone vanishes
identically on its tangent image.  Smoothness of a deformation base
and existence of a smooth fibre are distinct conclusions.  For a
degree-six point assigned $A_2$, a nonzero local vector alone does
not ensure a smoothing: each of its three scalar coefficients is
required by Theorem~\ref{thm04:criterion}.

If the singularities are only nodes and anticanonical cones over
$\dP_5$ or $\PP^1\times\PP^1$, the selected local bases are the entire
local bases.  Thus $\Def^S(X)=\Def(X)$ is smooth, of dimension
$h^1(X,\Theta_X)+\dim_\QQ K$.  Existence of a smoothing still
requires the nonzero relation of Theorem~\ref{thm04:criterion}.
